\section{Platooning}
\label{sec:sota_platooning}

Il \emph{platooning} consiste nella formazione e gestione di una colonna di veicoli che viaggiano con piccoli \emph{headway} mantenendo comportamento coordinato. Oltre ai benefici attesi in termini di fluidità del traffico ed efficienza, l’operatività del platoon richiede protocolli e strategie specifiche per supportare manovre come formazione, ingresso, uscita e gestione delle eccezioni, in modo coerente con i vincoli di sicurezza e con le capacità di comunicazione disponibili \cite{robinson2010operating,segata2014supporting}. La letteratura propone sia approcci centralizzati sia distribuiti, con meccanismi di consenso e coordinamento che permettono di validare e stabilizzare le manovre in scenari realistici, evidenziando come la componente di comunicazione intraplatoon sia strettamente intrecciata con la dinamica di controllo \cite{santini2017consensusbased,santini2019platooning}. In questa sezione si richiamano i concetti base e le principali modalità operative del platooning, ponendo l’attenzione sui requisiti che tali sistemi impongono a controllo e rete.
